% ═══════════════════════════════════════════════════════════════
% LEHRERHINWEISE: Waagerechter Wurf (Physik Klasse 10E)
% E-Kurs Gymnasium-Niveau (***)
% Doppelstunde: 70 Min. Erarbeitung + 20 Min. Präsentation
% ═══════════════════════════════════════════════════════════════

\documentclass[11pt, a4paper]{article}

% ═══════════════════════════════════════════════════════════════
% PAKETE
% ═══════════════════════════════════════════════════════════════

\usepackage[utf8]{inputenc}
\usepackage[T1]{fontenc}
\usepackage[ngerman]{babel}
\usepackage{helvet}
\renewcommand{\familydefault}{\sfdefault}

\usepackage[margin=2cm, top=2.5cm, bottom=2cm]{geometry}
\usepackage{enumitem}
\usepackage{tabularx}
\usepackage{booktabs}
\usepackage{array}
\usepackage{longtable}
\usepackage{multirow}
\usepackage{amsmath}
\usepackage{amssymb}

\usepackage{xcolor}
\usepackage{tcolorbox}
\usepackage{fancyhdr}
\usepackage{lastpage}

% Listen
\setlist{nosep, leftmargin=1.5em}

% Sternchen für Schwierigkeit
\newcommand{\stmark}{$\star$}

% ═══════════════════════════════════════════════════════════════
% FARBEN
% ═══════════════════════════════════════════════════════════════

\definecolor{physik}{HTML}{2c5aa0}
\definecolor{grau}{HTML}{888888}
\definecolor{hellgrau}{HTML}{f0f0f0}
\definecolor{success}{HTML}{2a7a4b}
\definecolor{warning}{HTML}{b35c00}
\definecolor{error}{HTML}{c62828}

\newcommand{\fachfarbe}{physik}

% ═══════════════════════════════════════════════════════════════
% VARIABLEN
% ═══════════════════════════════════════════════════════════════

\newcommand{\thema}{Waagerechter Wurf}
\newcommand{\fach}{Physik}
\newcommand{\klasse}{Klasse 10E (E-Kurs)}
\newcommand{\zeit}{90 Minuten (Doppelstunde)}

% ═══════════════════════════════════════════════════════════════
% KOPF- UND FUSSZEILE
% ═══════════════════════════════════════════════════════════════

\pagestyle{fancy}
\fancyhf{}
\renewcommand{\headrulewidth}{0.4pt}
\renewcommand{\footrulewidth}{0pt}

\lhead{\small\fach{} \klasse{} -- Lehrerhinweise}
\rhead{\small\thema}

\cfoot{\small\thepage/\pageref{LastPage}}

% ═══════════════════════════════════════════════════════════════
% BOXEN
% ═══════════════════════════════════════════════════════════════

\newtcolorbox{infobox}[1][]{
    colback=hellgrau,
    colframe=\fachfarbe,
    boxrule=0.8pt,
    arc=0pt,
    left=8pt, right=8pt, top=8pt, bottom=8pt,
    fonttitle=\bfseries,
    title=#1
}

\newtcolorbox{warnbox}[1][]{
    colback=white,
    colframe=warning,
    boxrule=1.5pt,
    arc=0pt,
    left=8pt, right=8pt, top=8pt, bottom=8pt,
    fonttitle=\bfseries,
    title=#1
}

\newtcolorbox{errorbox}[1][]{
    colback=white,
    colframe=error,
    boxrule=1.5pt,
    arc=0pt,
    left=8pt, right=8pt, top=8pt, bottom=8pt,
    fonttitle=\bfseries,
    title=#1
}

\newcommand{\abschnitt}[1]{%
    \vspace{12pt}
    \noindent\textbf{\large\textcolor{\fachfarbe}{#1}}
    \vspace{6pt}
    \hrule
    \vspace{8pt}
}

% ═══════════════════════════════════════════════════════════════
% DOKUMENT
% ═══════════════════════════════════════════════════════════════

\begin{document}

% ═══════════════════════════════════════════════════════════════
% TITEL
% ═══════════════════════════════════════════════════════════════

\begin{center}
{\LARGE\bfseries Lehrerhinweise}\\[4pt]
{\Large\textcolor{\fachfarbe}{\thema}}\\[8pt]
{\normalsize \klasse{} | \zeit{} | E-Kurs (\stmark\stmark\stmark-Niveau)}
\end{center}

\vspace{4pt}
\hrule height 1.5pt
\vspace{12pt}

% ═══════════════════════════════════════════════════════════════
% ÜBERSICHT
% ═══════════════════════════════════════════════════════════════

\abschnitt{Stundenübersicht}

\begin{infobox}[Kernziele der Stunde]
\begin{enumerate}
    \item Das Überlagerungsprinzip verstehen und anwenden
    \item Fallzeit und Wurfweite berechnen können
    \item \textbf{(E-Kurs):} Resultierende Geschwindigkeit $v_{res}$ und Auftreffwinkel $\alpha$ berechnen
    \item Physikalisches Wissen auf reale Situationen übertragen
\end{enumerate}
\end{infobox}

\vspace{8pt}

\begin{tabular}{llp{8cm}}
\textbf{Rahmenplan:} & \multicolumn{2}{l}{Kinematik -- Zusammengesetzte Bewegungen} \\
\textbf{Lernpfad:} & \multicolumn{2}{l}{LP-01-waagerechter-wurf-de.html (bzw. -en.html)} \\
\textbf{Simulation:} & \multicolumn{2}{l}{SIM-01-waagerechter-wurf.html} \\
\textbf{Arbeitsblatt:} & \multicolumn{2}{l}{AB-01-waagerechter-wurf-de.tex (bzw. -en.tex)} \\
\textbf{Musterlösung:} & \multicolumn{2}{l}{ML-01-waagerechter-wurf-de.tex (bzw. -en.tex)} \\
\end{tabular}

% ═══════════════════════════════════════════════════════════════
% STUNDENVERLAUF
% ═══════════════════════════════════════════════════════════════

\abschnitt{Stundenverlauf (Detailplanung)}

\renewcommand{\arraystretch}{1.3}
\begin{longtable}{|c|p{3.5cm}|p{6cm}|p{3cm}|}
\hline
\textbf{Zeit} & \textbf{Phase} & \textbf{Inhalt/Aktivität} & \textbf{Medium/Sozialform} \\
\hline
\endfirsthead
\hline
\textbf{Zeit} & \textbf{Phase} & \textbf{Inhalt/Aktivität} & \textbf{Medium/Sozialform} \\
\hline
\endhead
\hline
\endfoot

0--5' & Einstieg & Begrüßung, Ausgabe AB, Erklärung des Ablaufs & Plenum \\
\hline
5--10' & Problemstellung & LP Schritt 1: Was passiert? Hypothese aufstellen & EA mit Laptop/Tablet \\
\hline
10--20' & Simulation & LP Schritt 2: Beobachtung in der Simulation & EA mit Simulation \\
\hline
20--30' & Erklärung & LP Schritt 3: Überlagerungsprinzip erarbeiten, AB Kasten 1 ausfüllen & EA, AB \\
\hline
30--40' & Formeln I & LP Schritt 4: Grundformeln ($t$, $s_x$), AB Kasten 2 & EA, AB \\
\hline
40--50' & Formeln II & LP Schritt 5: Erweiterte Formeln ($v_y$, $v_{res}$, $\alpha$), AB Kasten 3 & EA, AB \\
\hline
50--55' & Relevanz & LP Schritt 6: ``Was hat das mit mir zu tun?'' & EA \\
\hline
55--70' & Präsentation & LP Schritt 7: Präsentationsaufgabe wählen und vorbereiten & EA/PA \\
\hline
\multicolumn{4}{|c|}{\textbf{--- Pause falls nötig (5 Min.) ---}} \\
\hline
70--90' & Präsentation & Schülerpräsentationen (je 3--5 Min.), Bewertung & Plenum \\
\hline
\end{longtable}

% ═══════════════════════════════════════════════════════════════
% MATERIALIEN
% ═══════════════════════════════════════════════════════════════

\abschnitt{Benötigte Materialien}

\begin{itemize}
    \item \textbf{Digital:} Laptops/Tablets für alle SuS mit Internetzugang
    \item \textbf{Print:} AB-01-waagerechter-wurf (1x pro SuS)
    \item \textbf{Lehrkraft:} Beamer zur Demonstration der Simulation
    \item \textbf{Optional:} Taschenrechner (für $\arctan$-Funktion)
    \item \textbf{Optional:} 2 identische Bälle für Live-Demonstration
\end{itemize}

% ═══════════════════════════════════════════════════════════════
% DIFFERENZIERUNG
% ═══════════════════════════════════════════════════════════════

\abschnitt{Differenzierung}

\begin{tabular}{|p{2.5cm}|p{5cm}|p{6cm}|}
\hline
\textbf{Niveau} & \textbf{Fokus} & \textbf{Unterstützung} \\
\hline
\stmark{} Basis & Überlagerungsprinzip verstehen, Grundformeln anwenden & Tabelle mit Fallzeiten nutzen, Aufgaben 1--2 \\
\hline
\stmark\stmark{} Standard & Alle Grundformeln beherrschen, Anwendungsaufgaben lösen & Aufgaben 1--3, Präsentation Option A oder B \\
\hline
\stmark\stmark\stmark{} E-Kurs & Erweiterte Formeln ($v_{res}$, $\alpha$), Transferaufgaben & Alle Aufgaben, anspruchsvolle Präsentation (Option C) \\
\hline
\end{tabular}

\vspace{8pt}

\begin{warnbox}[Hinweis zur Präsentationsbewertung]
Die drei Präsentationsoptionen sind nach Schwierigkeit gestaffelt:
\begin{itemize}
    \item \textbf{Option A (Klippenspringer):} Grundberechnung -- für alle machbar
    \item \textbf{Option B (Drohnenlieferung):} Anwendung mit Kontext -- mittleres Niveau
    \item \textbf{Option C (Stuntsprung):} Rückwärtsrechnung -- anspruchsvoll
\end{itemize}
Empfehlung: Leistungsstärkere SuS ermutigen, Option C zu wählen.
\end{warnbox}

% ═══════════════════════════════════════════════════════════════
% HÄUFIGE FEHLER
% ═══════════════════════════════════════════════════════════════

\abschnitt{Häufige Schülerfehler und Reaktionen}

\begin{errorbox}[Typische Fehler]
\begin{enumerate}
    \item \textbf{``Der schnellere Ball fliegt weiter, also fällt er länger.''}

    \textit{Reaktion:} Simulation nochmal zeigen. Beide Bälle starten gleichzeitig -- wann landen sie? Die horizontale Geschwindigkeit beeinflusst nur die Weite, nicht die Fallzeit!

    \item \textbf{Einheiten vergessen oder verwechseln}

    \textit{Reaktion:} ``Was passiert, wenn du m/s mit s multiplizierst? Die Sekunden kürzen sich weg -- was bleibt?''

    \item \textbf{$v_{res} = v_x + v_y$ (statt Pythagoras)}

    \textit{Reaktion:} Vektoraddition vs. Betrag. Skizze zeichnen lassen: Rechtwinkliges Dreieck! Bei senkrechten Vektoren \textbf{immer} Pythagoras.

    \item \textbf{Falsche Formel für Fallzeit: $t = \sqrt{h/g}$ (Faktor 2 vergessen)}

    \textit{Reaktion:} Herleitung aus $h = \frac{1}{2}gt^2$ zeigen. Die $\frac{1}{2}$ wird beim Umstellen zu 2 im Zähler.

    \item \textbf{Auftreffwinkel zur Vertikalen statt zur Horizontalen}

    \textit{Reaktion:} In der Skizze verdeutlichen: $\alpha$ ist der Winkel zur \textit{horizontalen} Achse (zum Boden).
\end{enumerate}
\end{errorbox}

% ═══════════════════════════════════════════════════════════════
% HINWEIS g-WERT
% ═══════════════════════════════════════════════════════════════

\abschnitt{Hinweis zur Erdbeschleunigung}

\begin{infobox}[Warum $g \approx 10\,\text{m/s}^2$?]
Wir verwenden $g \approx 10\,\text{m/s}^2$ statt $g = 9{,}81\,\text{m/s}^2$ aus didaktischen Gründen:

\begin{itemize}
    \item \textbf{Kopfrechenbar:} Schöne Zahlen ermöglichen das Rechnen ohne Taschenrechner
    \item \textbf{Fokus auf Konzept:} Weniger Rechenaufwand $\rightarrow$ mehr Zeit für Verständnis
    \item \textbf{Fehler < 2\%:} Der Näherungsfehler ist für Schulzwecke vernachlässigbar
\end{itemize}

\textbf{Im Lernpfad und AB ist dieser Hinweis enthalten.} Bei Nachfragen erklären, dass in der Praxis (Ingenieurwesen, Physik-Studium) mit exakten Werten gerechnet wird.
\end{infobox}

% ═══════════════════════════════════════════════════════════════
% BEWERTUNG PRÄSENTATION
% ═══════════════════════════════════════════════════════════════

\abschnitt{Bewertungskriterien Präsentation (20 Min. Phase)}

\begin{tabular}{|p{4cm}|c|p{7cm}|}
\hline
\textbf{Kriterium} & \textbf{Punkte} & \textbf{Beschreibung} \\
\hline
Korrekte Berechnung & 8 & Alle Formeln richtig angewendet, Rechenwege nachvollziehbar \\
\hline
Physikalische Erklärung & 4 & Zusammenhang zum Überlagerungsprinzip hergestellt \\
\hline
Darstellung & 4 & Skizze, übersichtliche Präsentation \\
\hline
Kommunikation & 4 & Verständliche Erklärung, Fachsprache \\
\hline
\textbf{Gesamt} & \textbf{20} & \\
\hline
\end{tabular}

\vspace{8pt}

\textbf{Umrechnung in Notenpunkte (15-NP-Skala):}

\begin{tabular}{|c|c|c|c|c|c|c|c|c|}
\hline
Punkte & 20--19 & 18--17 & 16--15 & 14--13 & 12--11 & 10--9 & 8--6 & <6 \\
\hline
NP & 15--14 & 13--12 & 11--10 & 9--8 & 7--6 & 5--4 & 3--1 & 0 \\
\hline
\end{tabular}

% ═══════════════════════════════════════════════════════════════
% KURZLÖSUNGEN
% ═══════════════════════════════════════════════════════════════

\abschnitt{Kurzlösungen (Referenz)}

\renewcommand{\arraystretch}{1.2}
\begin{tabular}{|l|l|l|}
\hline
\textbf{Aufgabe} & \textbf{Teilaufgabe} & \textbf{Lösung} \\
\hline
\multirow{2}{*}{1 (Ball, 20m, 15m/s)} & a) Fallzeit & $t = 2\,\text{s}$ \\
& b) Wurfweite & $s_x = 30\,\text{m}$ \\
\hline
\multirow{2}{*}{2 (Flugzeug, 45m, 20m/s)} & a) Flugzeit & $t = 3\,\text{s}$ \\
& b) Abwurfentfernung & $s_x = 60\,\text{m}$ \\
\hline
\multirow{4}{*}{3 (Klippenspringer, 80m, 30m/s)} & a) Fallzeit & $t = 4\,\text{s}$ \\
& b) Wurfweite & $s_x = 120\,\text{m}$ \\
& c) Vertikalgeschw. & $v_y = 40\,\text{m/s}$ \\
& d) Aufprallgeschw. & $v_{res} = 50\,\text{m/s}$ \\
\hline
\multirow{2}{*}{4 (Stuntman, 90m, 50m/s)} & a) Anfangsgeschw. & $v_0 = 30\,\text{m/s}$ \\
& b) Fallzeit & $t = 3\,\text{s}$ \\
\hline
\end{tabular}

\vspace{8pt}

\textbf{Präsentationsaufgaben:}

\begin{tabular}{|p{3.5cm}|p{3cm}|p{3cm}|p{3cm}|}
\hline
\textbf{Option} & \textbf{t} & \textbf{$s_x$} & \textbf{$v_{res}$} \\
\hline
A: Klippenspringer (45m, 25m/s) & 3 s & 75 m & $\approx 39\,\text{m/s}$ \\
\hline
B: Drohne (80m, 15m/s) & 4 s & 60 m & $\approx 43\,\text{m/s}$ \\
\hline
C: Stuntsprung ($s_x=100$m, 25m/s) & 4 s & (gegeben) & 50 m/s \\
\hline
\end{tabular}

% ═══════════════════════════════════════════════════════════════
% ERWEITERUNG
% ═══════════════════════════════════════════════════════════════

\abschnitt{Mögliche Erweiterungen / Hausaufgabe}

\begin{itemize}
    \item \textbf{Recherche:} Wie funktioniert ein Targetpod bei Kampfflugzeugen? (Ballistikcomputer)
    \item \textbf{Experiment:} Video-Analyse eines echten waagerechten Wurfs mit Smartphone
    \item \textbf{Vorbereitung nächste Stunde:} Was ändert sich beim \textit{schrägen} Wurf?
\end{itemize}

% ═══════════════════════════════════════════════════════════════
% FUSSZEILE
% ═══════════════════════════════════════════════════════════════

\vfill
\hrule
\vspace{4pt}
\footnotesize
\textit{Erstellt für Greenhouse Schools Rostock | E-Kurs Physik Klasse 10 | Bewegungslehre}

\end{document}
